% Incorporación propuesta; no incluida automáticamente en main.tex.
% Resultado recuperado y formalización reunida, 2026-09-10.
% Propietarios completos: registro_incidencias.tex y constructor_incidencial.py
% del Artículo II REV09; procedencia detallada en la nota Markdown homónima.
\subsection{Evaluación del libro orientado y reconstrucción integral}
\label{vi:lib:seccion}

La evaluación incidencial tiene un dominio distinto de su selección
dinámica. El dominio reúne visitas APP y trazas orientadas; la selección
TPK determina qué libro corresponde al estado terminal. Una vez fijado
ese libro, la evaluación siguiente es finita, unívoca y enteramente
aritmética. La prueba permite identificar exactamente qué información
se publica y qué información permanece en el historial.

\begin{definition}[Libro incidencial de doce ventanas]
\label{vi:lib:dominio}
Un libro \(\mathscr L\) consta de una familia finita no vacía de rutas,
un conjunto finito de visitas identificadas y una familia finita de
trazas orientadas. Cada ruta conserva al menos sus visitas de los
instantes \(1,\ldots,108\), con una sola visita por instante y ruta.
Se permiten visitas posteriores cuando una traza las necesite. Una
visita \(v\) conserva
\[
 (\gamma(v),t(v),i(v),j(v),s(v),\tau(v),w(v),\epsilon_\partial(v)),
\]
donde \(1\le i,j\le9\), \(\tau\in\{-1,0,1\}\),
\(w\) es una multiplicidad entera positiva y \(s\) designa la hoja.
La lectura aritmética declarada en esta carta es
\[
 n_v=\begin{cases}
 i(v)+j(v),&s(v)=\Sigma,\\
 i(v)j(v),&s(v)=\Pi,\\
 9,&s(v)=0,
 \end{cases}
 \qquad n_v=r_v+9q_v,\quad1\le r_v\le9.
\]
La hoja \(0\) es la convención de umbral de esta carta. Para una visita
con \(r_v=9\), el operador de frontera especifica
\(\epsilon_\partial(v)=1\) en la corona y \(-1\) en el interior;
esta etiqueta no se deduce únicamente del valor del residuo.

Cada traza \(\theta=(v_1,\ldots,v_N)\) está contenida en una ruta,
sus instantes son consecutivos y su origen pertenece a una ventana
\(E_m=\{9m-8,\ldots,9m\}\). Conserva una multiplicidad positiva
\(u(\theta)\) y una unidad de longitud reducida \(\lambda(\theta)>0\)
entera, de modo que
\[
 N\lambda(\theta)=120(1+3j(\theta)),\qquad j(\theta)\ge0.
\]
Su canal orientado es
\[
 \sigma(\theta)=\operatorname{sgn}
       \left(\sum_{a=1}^{N}\tau(v_a)\right)\in\{-1,1\}.
\]
Esta familia sólo cuenta trazas de canal no nulo. La selección de la
familia, las multiplicidades, el operador de frontera y las unidades
de longitud son datos declarados del libro. Si el libro se presenta
como publicación del TPK, su procedencia incluye además la regla que
los obtiene del transporte y la cobertura de las trazas admitidas por
esa regla. La comprobación de formato de un libro no demuestra esa
procedencia.
\end{definition}

La división positiva anterior se calcula, sin recibir el residuo como
otro argumento, mediante
\[
 r_v=1+((n_v-1)\bmod9),\qquad q_v=(n_v-r_v)/9.
\]
La diferencia entre longitud reducida y número de instantes se conserva
en \(\lambda\); el calendario de \(108\) instantes no autoriza a
identificar ambas magnitudes por omisión.

Sea \(\mathscr V_m\) el conjunto de visitas cuyo instante está en
\(E_m\), y \(\mathscr T_m\) la familia de trazas cuyo origen está en
esa ventana. Se calculan los tres recuentos
\begin{align}
 A_m&=\sum_{v\in\mathscr V_m}w(v)
              \mathbf1_{\{1,3,5,7\}}(r_v),\notag\\
 C_m&=-\sum_{\theta\in\mathscr T_m}u(\theta)\sigma(\theta),\notag\\
 V_m&=\sum_{v\in\mathscr V_m}w(v)
              \mathbf1_{\{9\}}(r_v)\epsilon_\partial(v).
 \label{vi:lib:recuentos}
\end{align}
El segundo coincide con
\(\sum_j(\nu^-_{120,m}(j)-\nu^+_{120,m}(j))\), donde
\(\nu^\pm\) es la suma de multiplicidades de las trazas del canal
correspondiente y longitud \(120(1+3j)\). Todas las sumas son finitas.
Para cualquier entero \(x\), sea
\(x=[x]_{10}+10\lfloor x/10\rfloor\), con \(0\le[x]_{10}<10\).
Los bloques y sus coordenadas transversales se definen por
\begin{align}
 z_m&=100[A_m]_{10}+10[C_m]_{10}+[V_m]_{10},\notag\\
 \mathcal E_{\mathrm{cnt}}(\mathscr L)
   &=\mathcal D(z):=(D_3z,D_4z,Q(z)),\label{vi:lib:directa}\\
 (D_hz)_i&=z_i-z_{i+h},\qquad Q(z)=\sum_i z_i.\notag
\end{align}
En las fórmulas cíclicas siguientes se numeran las posiciones de cero
a once, sin cambiar su orden. La aplicación tiene \(25\) coordenadas
enteras, pero su imagen tiene rango \(12\): las diferencias y la suma
no constituyen \(25\) datos independientes.

\begin{theorem}[Composición y reversibilidad sobre el libro declarado]
\label{vi:lib:composicion}
Para todo libro de la definición~\ref{vi:lib:dominio}, la fórmula
\eqref{vi:lib:directa} produce un único elemento de
\(\mathcal D(\{0,\ldots,999\}^{12})\). Su reconstrucción integral
es única. El cambio de representación armónico descrito abajo satisface
\[
 \Pi_H\mathcal E_{\mathrm{cnt}}(\mathscr L)=H_{12}z,
 \qquad
 \tfrac14H_{12}\Pi_H\mathcal E_{\mathrm{cnt}}(\mathscr L)=z.
\]
En particular, el sello obtenido por esta composición es \(K(\mathscr L)=z\).
Ninguna congruencia de compatibilidad debe ajustarse después de conocer
el resultado: todas son consecuencias de la evaluación directa.
\end{theorem}

\begin{proof}
Las sumas enteras y las divisiones euclídeas determinan unívocamente
los bloques. Sean \(b=D_3z\), \(c=D_4z\), \(q=Q(z)\), y
\[
 w_i=c_i-b_{i+1},\quad P_0=0,\quad
 P_i=\sum_{j<i}w_j,\quad T=\sum_{i=0}^{11}P_i.
\]
Sustituir la definición de las diferencias da
\(w_i=z_i-z_{i+1}\). Por tanto,
\[
 \sum_iw_i=0,\quad b_i=w_i+w_{i+1}+w_{i+2},\quad
 q+T=12z_0.
\]
Se obtiene la reconstrucción
\[
 z_i=(q+T)/12-P_i.
\]
Estas tres relaciones caracterizan también la imagen integral de
\(\mathcal D\). En efecto, si se cumplen para \(b,c,q\) y
\(q+T\equiv0\pmod{12}\), la fórmula reconstruida es entera,
cíclica, tiene diferencias \(b,c\) y suma \(q\). Su diferencia
con otra solución tiene todas las diferencias consecutivas nulas y
suma cero, por lo que es nula. La linealidad y esta inyectividad de
\(\mathcal D\colon\mathbb Z^{12}\to\mathbb Z^{25}\) prueban
la afirmación de rango.

Para escribir el cambio armónico, sean \(r=0,1,2\),
\[
 S_r=\sum_{j=0}^3c_{r+3j},\quad
 a_0=(q+2S_0+S_1)/3,\quad
 a_1=a_0-S_0,\quad a_2=a_1-S_1,
 \qquad d_{r,j}=b_{r+3j}.
\]
Definimos sobre la imagen integral, en el orden de cada órbita,
\[
 \Pi_H^{(r)}(b,c,q)=
 (a_r,d_{r,1}-d_{r,3},d_{r,0}+d_{r,2},d_{r,0}-d_{r,2}).
\]
Si \(s_r=\sum_jz_{r+3j}\), entonces
\(S_r=s_r-s_{r+1}\) y \(q=s_0+s_1+s_2\). De aquí
\(a_r=s_r\); en particular, sus divisiones por tres son enteras.
En la órbita \(x_j=z_{r+3j}\), se tiene \(d_{r,j}=x_j-x_{j+1}\).
La sustitución en las tres diferencias anteriores da
\[
 \Pi_H^{(r)}(b,c,q)=
 \begin{pmatrix}
 1&1&1&1\\1&1&-1&-1\\
 1&-1&1&-1\\1&-1&-1&1
 \end{pmatrix}
 \begin{pmatrix}x_0\\x_1\\x_2\\x_3\end{pmatrix}.
\]
Llamando \(H_4\) a esa matriz, \(H_4^2=4I_4\) por
multiplicación directa. \(H_{12}\) aplica estos tres bloques a las
órbitas ordenadas \((0,3,6,9)\), \((1,4,7,10)\),
\((2,5,8,11)\), conservando la misma convención al volver a las
doce posiciones. Por tanto \(H_{12}^2=4I_{12}\) y ambas
composiciones anunciadas se siguen. La preimagen empleada en esta
prueba no se introduce como argumento de \(\Pi_H\).
\end{proof}

\begin{corollary}[Información conservada por la publicación]
\label{vi:lib:fibras}
Sean \(\mathscr L,\mathscr L'\) dos libros y \(N,N'\) sus
matrices de recuentos \((A,C,V)\). Entonces
\[
 \mathcal E_{\mathrm{cnt}}(\mathscr L)
   =\mathcal E_{\mathrm{cnt}}(\mathscr L')
 \quad\Longleftrightarrow\quad
 N-N'\in10\mathbb Z^{36}.
\]
La salida conserva los \(36\) residuos decimales. Los cocientes de
división, la identidad de las visitas, sus hojas, sus orientaciones y
su transporte permanecen en el libro, pero no se recuperan en general
de esas \(25\) coordenadas.
\end{corollary}
\begin{proof}
La igualdad de los residuos produce los mismos bloques y coordenadas.
Recíprocamente, la inyectividad de \(\mathcal D\) fuerza \(z=z'\).
La unicidad de la escritura de cada entero entre \(0\) y \(999\)
en tres cifras recupera los tres residuos de cada ventana. La división
euclídea muestra que esto equivale a la congruencia indicada.
\end{proof}

La congruencia describe las fibras de una aplicación de conjuntos.
La codificación mediante representantes decimales no es aquí un
homomorfismo aditivo hacia \(\mathbb Z^{25}\); por ello no se
interpreta la afirmación como un núcleo lineal de dimensión \(36\).

\paragraph{Empalme con la selección dinámica.}
Sea \(\mathcal R_{12}(x)\) la extracción del historial de un estado
TPK, y sea \(\mathcal B\) la regla de selección que, de ese historial,
obtiene la familia incidencial de la definición~\ref{vi:lib:dominio}.
Cuando se especifica esa regla, la composición efectiva es
\[
 x\longmapsto\mathcal R_{12}(x)
  \longmapsto\mathcal B(\mathcal R_{12}(x))
  \longmapsto\mathcal E_{\mathrm{cnt}}
       (\mathcal B(\mathcal R_{12}(x)))
  \longmapsto U(x)\longmapsto K(x).
\]
La igualdad de su tramo incidencial con
\(\mathcal E_{108}^{90,120}\) significa que se usa el libro y el
lector del estado terminal correspondiente. El teorema demuestra todo
el tramo desde el libro suministrado; identificar ese libro con el
canónico requiere la selección \(\mathcal B\), sus multiplicidades,
su clasificación de frontera y su cobertura de trazas. Fijar únicamente
el calendario, la semilla inicial o un vector de doce eventos no
enumera por sí solo estos datos. Tampoco la inversión del sello
selecciona retrospectivamente el libro.

Así queda cerrado el resultado relativo de evaluación y reversibilidad
con estructura de partida explícita. La afirmación de que un estado
terminal concreto produce un sello numérico concreto añade una
evaluación de \(\mathcal B\mathcal R_{12}\) que debe conservarse
junto a esa afirmación; no se deduce de la invertibilidad recién
demostrada.
